\documentclass[8pt,twocolumn]{article}
\usepackage[spanish]{babel}
\usepackage[utf8]{inputenc}
\usepackage[margin=0.35in]{geometry}
\usepackage{amsmath,amssymb}
\usepackage{titlesec}

\titleformat{\section}{\small\bfseries\uppercase}{}{0em}{}[\titlerule]
\titleformat{\subsection}{\footnotesize\bfseries}{}{0em}{}
\titlespacing*{\section}{0pt}{4pt}{2pt}
\titlespacing*{\subsection}{0pt}{3pt}{1pt}
\setlength{\parindent}{0pt}
\setlength{\parskip}{2pt}
\pagestyle{empty}

\begin{document}

\begin{center}
    \textbf{\large FORMULARIO DE SUPERVIVENCIA: CC1002 -- CONTROL 1} \\
    \footnotesize Felipe Colli Olea -- FCFM Universidad de Chile
\end{center}

\section{1. La Receta de Diseño (Puntaje Seguro)}
\textbf{Obligatorio en CADA función para no perder hasta 1.5 pts:}
\begin{verbatim}
# nombreFuncion: tipo1 tipo2 -> tipoRetorno
# Proposito: Breve descripcion de que calcula
# Ejemplo: nombreFuncion(arg1, arg2) devuelve resultado
def nombreFuncion(param1, param2):
    # Cuerpo de la funcion
    return resultado

# Tests (Al menos 1 o 2 por funcion):
assert nombreFuncion(val1, val2) == esperado
assert not funcionBooleana(casoFalso)
assert cerca(funcionReal(val), esperado, 0.0001)
\end{verbatim}
\textbf{Tipos para contratos:} \texttt{int}, \texttt{float}, \texttt{num} (ambos), \texttt{str}, \texttt{bool}, \texttt{None} (procedimientos).

\section{2. Manipulación de Dígitos (Machete Clave)}
\textbf{Dado un número entero positivo $N$:}
\begin{itemize}
    \item \textbf{Último dígito (derecha):} \texttt{n \% 10}
    \item \textbf{Quitar último dígito:} \texttt{n // 10}
    \item \textbf{Último par de dígitos (base 100):} \texttt{n \% 100}
    \item \textbf{Quitar último par de dígitos:} \texttt{n // 100}
    \item \textbf{Dígito en posición $k$ (desde der=0):} \texttt{(n // (10**k)) \% 10}
    \item \textbf{Primeros 2 dígitos (izq):} \texttt{n // (10 ** (digitos(n) - 2))}
    \item \textbf{Quitar primeros 2 dígitos (izq):} \texttt{n \% (10 ** (digitos(n) - 2))}
    \item \textbf{Quitar ambos extremos (2 dígitos c/u):} \\
    \texttt{nuevo = (n \% (10**(largo - 2))) // 100}
\end{itemize}

\section{3. Plantillas de Recursión Clásicas}

\subsection{Contar Dígitos / Largo}
\begin{verbatim}
# digitos: int -> int
def digitos(n):
    if n < 10:
        return 1
    return 1 + digitos(n // 10)
\end{verbatim}

\subsection{Suma de Dígitos}
\begin{verbatim}
# sumaDigitos: int -> int
def sumaDigitos(n):
    if n < 10:
        return n
    return (n % 10) + sumaDigitos(n // 10)
\end{verbatim}

\subsection{Invertir un Número}
\begin{verbatim}
# invertir: int -> int
def invertir(n):
    if n < 10:
        return n
    k = digitos(n // 10)
    return (n % 10) * (10 ** k) + invertir(n // 10)
\end{verbatim}

\subsection{Mayor Dígito de un Número}
\begin{verbatim}
# maxDigito: int -> int
def maxDigito(n):
    if n < 10:
        return n
    resto = maxDigito(n // 10)
    u = n % 10
    return u if u > resto else resto
\end{verbatim}

\subsection{Procesar Secuencia por Bloques de 2 Dígitos}
\begin{verbatim}
# sumarPares: int -> int
def sumarPares(secuencia):
    if secuencia == 0:
        return 0
    return (secuencia % 100) + sumarPares(secuencia // 100)
\end{verbatim}

\subsection{Procedimiento Recursivo (Imprimir sin retornar)}
\begin{verbatim}
# mostrarPares: int -> None
def mostrarPares(n):
    if n > 0:
        print("Valor:", n % 100)
        mostrarPares(n // 100)
\end{verbatim}

\subsection{Comprobar Simetría de Extremos (Estilo Sumeru)}
\begin{verbatim}
# comprobar: int -> bool
def comprobar(N):
    largo = digitos(N)
    if largo == 4:
        return True
    # Extraer extremos
    der = N % 100
    izq = N // (10 ** (largo - 2))
    # Reducir número eliminando ambos extremos
    nuevoN = (N % (10 ** (largo - 2))) // 100
    if (izq + der) != sumarExtremos(nuevoN):
        return False
    return comprobar(nuevoN)
\end{verbatim}

\subsection{Función Wrapper con Acumulador (Recursión de Cola)}
\begin{verbatim}
# helper: int int -> int
def helper(n, acc):
    if n == 0:
        return acc
    return helper(n // 10, acc + (n % 10))

# funcionPrincipal: int -> int
def funcionPrincipal(n):
    return helper(n, 0)
\end{verbatim}

\section{4. Comparación de Flotantes (Cerca)}
\textbf{¡Nunca usar \texttt{==} para números reales (IEEE 754)!}
\begin{verbatim}
# cerca: num num num -> bool
def cerca(x, y, eps=0.0001):
    return abs(x - y) < eps

assert cerca(0.1 + 0.2, 0.3)
\end{verbatim}

\section{5. Módulos y Entradas/Salidas}
\subsection{Módulo Math y Random}
\begin{verbatim}
import math, random
math.sqrt(x)    # Raiz cuadrada
math.pi         # Constante pi
round(num, n)   # Redondea a n decimales
random.randint(a, b)  # Entero aleatorio en [a, b]
random.random()       # Float aleatorio en [0.0, 1.0)
\end{verbatim}

\subsection{Interacción Básica}
\begin{verbatim}
x = int(input("Ingrese entero: "))
y = float(input("Ingrese real: "))
# Concatenacion limpia sin errores de tipo:
print("Resultado: " + str(x) + " y " + str(y))
\end{verbatim}

\section{6. Trucos y Errores Típicos a Evitar}
\begin{itemize}
    \item \textbf{NOTA DEL ENUNCIADO:} En casi todos los controles prohíben explícitamente usar \texttt{str()} para procesar números en las preguntas de recursión. \textbf{Usa solo \texttt{\% 10}, \texttt{// 10} y potencias de 10}.
    \item \textbf{Condicionales directos:} En vez de \texttt{if x \% 2 == 0: return True else: return False}, escribe directo: \texttt{return (x \% 2 == 0)}.
    \item \textbf{Importar módulos de la pregunta:} Si el enunciado dice que existe un módulo \texttt{domino} o \texttt{musical}, no olvides escribir al inicio: \\
    \texttt{from domino import *}
    \item \textbf{Pre-condición con assert:} Si te piden que $x > 0$, mete un \texttt{assert x > 0} dentro del cuerpo antes de la lógica.
\end{itemize}

\end{document}
